\subsection{Estructura Organizacional}

La estructura orgánica de la Universidad Nacional de Ingeniería está normada en su Reglamento de Organización y Funciones \cite{rof2023}. Para efectos del procedimiento de matrícula regular, intervienen de forma coordinada órganos de la alta dirección, dependencias centrales de registro y soporte, y las unidades académicas de la Facultad de Ciencias.

\paragraph{Nivel de Alta Dirección}
\begin{itemize}
    \item \textbf{Rectorado:} Ejerce la máxima representación legal y ejecutiva de la institución, conduciendo la gestión universitaria conforme a las políticas del Consejo Universitario.
    \item \textbf{Vicerrectorado Académico (VRAD):} Órgano de alta dirección responsable de normar, supervisar y evaluar las actividades académicas de pregrado y posgrado, incluyendo el calendario oficial de matrícula.
\end{itemize}

\paragraph{Órganos Centrales de Registro y Soporte}
\begin{itemize}
    \item \textbf{Dirección de Registros Centrales y Estadística (DIRCE / ORCE):} Administra la base de datos académica general y la plataforma web \texttt{INTRALU}. Es la encargada de habilitar las cuentas de usuario, publicar el cronograma por turno ponderado y consolidar la matrícula oficial a nivel universitario.
    \item \textbf{Centro Médico y Dirección de Bienestar Universitario:} Gestiona el Fondo de Autoseguro Estudiantil, atiende las emergencias médicas y registra las validaciones de salud requeridas para la matrícula.
    \item \textbf{Dirección General de Administración (DIGA):} Verifica el estado de pagos, aranceles complementarios o posibles deudas del estudiante con la universidad.
\end{itemize}

\paragraph{Nivel Descentralizado: Facultad de Ciencias}
\begin{itemize}
    \item \textbf{Decanato y Consejo de Facultad:} Máxima autoridad de gobierno de la facultad encargada de aprobar la programación curricular y los criterios específicos de matrícula de sus carreras.
    \item \textbf{Comisión de Matrícula:} Órgano especializado integrado por docentes y delegados estudiantiles que evalúa solicitudes de rectificación, ampliación de créditos y casos académicos especiales.
    \item \textbf{Escuelas Profesionales y Departamentos Académicos:} Diseñan los planes de estudio, definen las vacantes por asignatura y asignan a los docentes encargados de las secciones de teoría y laboratorio.
    \item \textbf{Oficina de Estadística y Registros Académicos (AERA):} Unidad técnica de la facultad que ingresa las secciones, cursos y horarios en el sistema antes de iniciar el ciclo.
\end{itemize}

En la Figura~\ref{fig:organigrama_uni} se presenta la estructura jerárquica y funcional de estas áreas vinculadas al proceso de matrícula.

% ==============================================================================
% Diagrama: Organigrama Funcional Institucional (Enfoque Proceso de Matrícula)
% Estilo modular y declarativo con macros reutilizables y posicionamiento relativo
% ==============================================================================

\begin{figure}[H]
\centering

% Definición de macros para la estructura del organigrama
\providecommand{\nodoOrg}[4][]{%
    % #1: opciones de posición relativa (ej. below=of..., right=of...)
    % #2: nombre o id del nodo
    % #3: texto de contenido
    % #4: estilo visual
    \node[#4, #1] (#2) {#3};
}

\providecommand{\conectarOrtogonal}[2]{%
    % Conexión estándar padre-hijo (salida por el sur, entrada ortogonal por el norte)
    \draw[lineaOrg] (#1.south) -- ++(0,-0.45) -| (#2.north);
}

\providecommand{\conectarLista}[2]{%
    % Conexión vertical tipo lista (salida por el oeste/sur hacia la izquierda del hijo)
    \draw[lineaOrg] (#1.west) -| (#2.west);
}

\begin{tikzpicture}[
    >=Latex,
    node distance=0.6cm and 0.6cm,
    % Estilos visuales por jerarquía funcional
    baseBox/.style={
        rectangle, 
        rounded corners=2pt, 
        draw=black!75, 
        thick, 
        align=center, 
        font=\small\sffamily,
        inner sep=5pt
    },
    nivelRector/.style={
        baseBox, 
        fill=blue!15, 
        text width=6.2cm, 
        minimum height=0.9cm, 
        font=\small\bfseries\sffamily
    },
    nivelCentral/.style={
        baseBox, 
        fill=cyan!12, 
        text width=3.8cm, 
        minimum height=0.85cm
    },
    nivelFacultad/.style={
        baseBox, 
        fill=orange!15, 
        text width=4.8cm, 
        minimum height=0.85cm, 
        font=\small\bfseries\sffamily
    },
    nivelUnidad/.style={
        baseBox, 
        fill=gray!10, 
        draw=black!60, 
        text width=3.8cm, 
        minimum height=0.75cm, 
        font=\footnotesize\sffamily
    },
    lineaOrg/.style={
        draw=black!70, 
        thick
    }
]

    % --------------------------------------------------------------------------
    % 1. Nivel Superior / Alta Dirección
    % --------------------------------------------------------------------------
    \nodoOrg{rectorado}{Rectorado / Consejo Universitario}{nivelRector}
    \nodoOrg[below=0.6cm of rectorado]{vrad}{Vicerrectorado Académico}{nivelRector}
    \draw[lineaOrg] (rectorado.south) -- (vrad.north);

    % Punto de bifurcación central debajo del VRAD
    \coordinate (bifurcacion) at ($(vrad.south) + (0,-0.45)$);
    \draw[lineaOrg] (vrad.south) -- (bifurcacion);

    % --------------------------------------------------------------------------
    % 2. Órganos Centrales de Apoyo y Registro (Lado Izquierdo)
    % --------------------------------------------------------------------------
    \nodoOrg[below left=0.8cm and 0.4cm of vrad]{dirce}{\textbf{DIRCE / ORCE}\\(Registros Centrales y Matrícula)}{nivelCentral}
    \nodoOrg[below=0.5cm of dirce]{salud}{\textbf{Centro Médico / Bienestar}\\(Autoseguro Estudiantil)}{nivelCentral}
    \nodoOrg[below=0.5cm of salud]{diga}{\textbf{DIGA / Tesorería}\\(Validación de Tasas / Deudas)}{nivelCentral}

    % --------------------------------------------------------------------------
    % 3. Órgano de Línea Académica Descentralizada: Facultad de Ciencias (Lado Derecho)
    % --------------------------------------------------------------------------
    \nodoOrg[below right=0.8cm and 0.4cm of vrad]{decanato}{Facultad de Ciencias\\(Decanato / Consejo de Facultad)}{nivelFacultad}
    
    % Unidades internas de la Facultad de Ciencias vinculadas a la matrícula
    \nodoOrg[below=0.5cm of decanato]{comision}{\textbf{Comisión de Matrícula}\\(Casos especiales y rectificaciones)}{nivelUnidad}
    \nodoOrg[below=0.4cm of comision]{escuelas}{\textbf{Escuelas Profesionales}\\(CC, Física, Matemática, Química)}{nivelUnidad}
    \nodoOrg[below=0.4cm of escuelas]{dptos}{\textbf{Departamentos Académicos}\\(Carga docente y horarios)}{nivelUnidad}
    \nodoOrg[below=0.4cm of dptos]{estadistica}{\textbf{Oficina de Estadística / AERA}\\(Apertura de cursos en sistema)}{nivelUnidad}

    % --------------------------------------------------------------------------
    % 4. Enlaces y ramificación
    % --------------------------------------------------------------------------
    % Conexión desde el punto de bifurcación a DIRCE y Decanato
    \draw[lineaOrg] (bifurcacion) -| (dirce.north);
    \draw[lineaOrg] (bifurcacion) -| (decanato.north);

    % Conexiones verticales en bloque central
    \draw[lineaOrg] (dirce.south) -- (salud.north);
    \draw[lineaOrg] (salud.south) -- (diga.north);

    % Conexiones verticales en bloque de facultad
    \draw[lineaOrg] (decanato.south) -- (comision.north);
    \draw[lineaOrg] (comision.south) -- (escuelas.north);
    \draw[lineaOrg] (escuelas.south) -- (dptos.north);
    \draw[lineaOrg] (dptos.south) -- (estadistica.north);

    % Agrupación visual por áreas
    \node[draw=blue!40, dashed, rounded corners, fit=(dirce)(salud)(diga), inner sep=6pt, label=above:{\footnotesize\bfseries\color{blue!80!black} Órganos Centrales de Soporte}] {};
    \node[draw=orange!50, dashed, rounded corners, fit=(decanato)(comision)(escuelas)(dptos)(estadistica), inner sep=6pt, label=above:{\footnotesize\bfseries\color{orange!80!black} Nivel de Facultad}] {};

\end{tikzpicture}
\caption{Estructura organizacional funcional vinculada al proceso de matrícula regular.}
\label{fig:organigrama_uni}
\end{figure}


